%%=============================================================================
%% Voorwoord
%%=============================================================================

\chapter*{\IfLanguageName{dutch}{Woord vooraf}{Preface}}%
\label{ch:voorwoord}

%% TODO:
%% Het voorwoord is het enige deel van de bachelorproef waar je vanuit je
%% eigen standpunt (``ik-vorm'') mag schrijven. Je kan hier bv. motiveren
%% waarom jij het onderwerp wil bespreken.
%% Vergeet ook niet te bedanken wie je geholpen/gesteund/... heeft

%%=============================================================================
%% Voorwoord
%%=============================================================================


My name is Joris Yangsheng Xu and I am a final year student in the Applied Information Technology programme at the University of Applied Sciences Ghent. 

This bachelor thesis was written as part of the \"bachelorproef\" subject within the Applied IT curriculum.
From the beginning of my studies, I was less focused on following a fixed path and more interested in learning technically challenging and valuable topics. 
I have always been naturally curious and often explored technologies outside of the standard curriculum. 
During my work based learning experience at Rabobank, this curiosity eventually led me towards IBM MQ and middleware technologies.

One of the moments that strongly influenced the direction of this thesis started rather unexpectedly. 
While independently studying IBM MQ during work based learning, I happened to sit next to Dragos, an experienced IBM MQ administrator. 
After noticing what I was studying, he spent a long time explaining IBM MQ concepts and internals in significant detail. 
At that moment, much of the explanation was still far beyond my level of understanding, but it introduced me to a completely new technical world. 
Instead of discouraging me, it motivated me to continue learning and asking questions.

What started as curiosity gradually evolved into the topic of this bachelor thesis. 
Through discussions with Dragos a retired former middleware engineer and by exploring my own ideas, I became interested in understanding a very specific and practical part of IBM MQ:\@ buffer pools. 
I wanted to become an owner of a small but important component within a real enterprise system and explore how operational data could be transformed into useful insights.

Working on this thesis taught me much more than only technical knowledge. 
Setting up environments, understanding certificates, working with USS, HTTP servers, SMF, log streams, parsing data, 
scheduling workloads, building APIs, and visualising data in Grafana repeatedly introduced new challenges that required additional learning. 
Each solved problem introduced several new things I did not yet understand. 
Looking back, one of the most valuable lessons was not simply learning new technologies, 
but developing a clearer understanding of the complexity of enterprise systems and becoming more aware of what I did not know.

Although the process was challenging at times, I do not consider it a negative experience. Strict deadlines and accountability from 
mentors pushed me further than I initially expected.  
The experience taught me the importance of persistence, communication, asking questions, and continuously improving through feedback.

Finally, I would like to thank several people who played an important role during this thesis.

First, I would like to thank Leendert for his guidance and support throughout the thesis. 
His help with structuring the thesis, limiting the scope when necessary, discussing progress, 
and calming nerves during stressful moments was invaluable.

I would also like to thank Dragos for his enthusiasm, sharing his technical expertise and experience in IBM MQ\@ His direct and honest feedback, technical discussions, and ability to call me out on my thinking greatly contributed to both the thesis and my personal growth.

In addition, I would like to thank Jordi for helping me structure the data pipeline and architecture recommendations, providing practical advice from his 
own experiences, and advocating for the time and support needed to complete this work.

Finally, I would like to thank everyone who contributed, directly or indirectly, to this learning experience and supported me throughout the process.